%!TEX root = ../thesis_main.tex
%!TEX encoding = UTF-8 Unicode

\section{Definition of tars in biomass gasification}
\label{sec:theo_tar}

Tar from biomass gasification is usually defined as follows: \textit{all organics boiling at temperatures above that of benzene should be considered as tar} \cite{Milne1998}. As it is not clear whether benzene itself is included in this definition, it will be treated separately and explicitly included or excluded in this study.

Tar species are usually classified as primary, secondary and tertiary tar, following the definition given in the review of \citet{Milne1998}. Primary tars are oxygenated compounds resulting from the pyrolysis of cellulose, hemicellulose and lignin, such as levoglucosan, glycolaldehyde, guaiacol and furfurals. Secondary tars are composed of a diversity of phenolics and olefins produced by cracking and oxidation of primary tars in mild conditions (500-800°C). Tertiary tars result from reactions of primary and secondary tars at higher temperature (above 800°C). They are composed of substituted (methyl derivatives) and unsubstituted aromatics and polyaromatic hydrocarbons (PAHs): benzene, toluene, naphthalene, indene, methylnaphthalene and PAHs with 3 or more rings. Primary and secondary tars can be thermally cracked into \ce{CO}, \ce{H2} and light gases, while tertiary tars species are refractory and much more difficult to destroy.

A complementary classification of tar species was proposed by \citet{Kiel2004}, consisting of 5 classes based on tar species condensation behavior and water solubility:
\begin{itemize}
\item Class I: GC-undetectable. Heaviest molecules that condense at very high temperature.
\item Class II: Heterocyclic compounds containing \ce{O} and/or \ce{N} atoms, mainly phenolics or nitrogenous aromatics. They condense at moderate temperatures, between 50 and 100°C for concentrations from 10 to 1000\,mg/Nm$^3$. Their polarity induces a high water solubility, increasing the complexity of waste-water treatment from the cleaning system.
\item Class III: Aromatics derivatives of benzene, such as toluene and xylene. They are not polar and do not condense thanks to their very low dew point (below 0°C up to 10\,g/Nm$^3$).
\item Class IV: Light PAHs of 2-3 rings, from naphthalene to phenanthrene. Non-polar, they condense at moderate temperatures similar to Class II species.
\item Class V: heavy PAHs of 4-7 rings, from fluoranthene to coronene. They condense at high temperature even when in very low concentration (above 100°C at 0.1\,mg/Nm$^3$).
\end{itemize}

\section{Mechanisms of \textit{in-situ} tar reduction}

The mechanism of tar reduction and destruction inside a gasifier can be roughly split in two main parts: the homogeneous conversion of primary and secondary tars into light gases and tertiary tars (aromatics and Polyaromatic Hydrocarbons -- PAHs), and the decomposition of tertiary tars by heterogeneous reactions over the char bed acting as a catalyst.  This is particularly applicable in a two-stage gasifier, where the two parts occur, respectively, within the oxidation and the reduction zones. The destruction of tar inside a gasifier is favoured by three factors: temperature, oxygen and steam, in that order of importance. High temperature partial oxidation is essential for efficient tar reforming, especially for the conversion of primary and secondary tars \cite{Zhang2010}. 

Steam can affect both parts of that process. By releasing \ce{H^.}, \ce{O^{..}} and \ce{OH^.} radicals at high temperature, it favors the conversion of substituted aromatics and phenols into benzene and light PAHs, and hence reduces the yield of heavy PAHs of which phenol is the main precursor. In particular, the presence of many \ce{H^.} radicals help to saturate intermediate radicals to single ring aromatics (benzene) and light hydrocarbons (\ce{C2}, \ce{C3}), avoiding their recombination into polyaromatic compounds \cite{Feng2017}.

Steam can also indirectly affect the heterogeneous conversion of heavy PAHs by maintaining the catalytic activity of char, through the gasification of char and coke deposits by \ce{OH^.} radicals. Under a sufficient temperature (700--900°C), tar heterogenous conversion occurs through 1) adsorption on the char active sites, 2) polymerization (coke formation) and 3) gasification.  This efficiently decomposes aromatics and heavy tar, but char activity rapidly decreases due to micropores blocking by coke deposition: steam can help maintain char activity by supporting the gasification step \cite{Hosokai2011}.

%Internal note: According to THUNMAN, the impact of steam is explained by the fact that it provides a lot of H+ radicals. The presence of these radicals helps to saturate aromatic rings to benzene, thereby avoiding their recombination in polyaromatic compounds. His explanation is as follows:
%o	Steam breaks down into H+ and OH- radicals.
%o	Sugar compounds break open by releasing the O in the ring, forming an unsaturated chain.
%o	If many H+ radicals are available, they saturate the chain in either light HCs (C2, C3) or single ring aromatics (benzene).
%o	The OH- radicals on the other hand react with carbon (char) to form CO2 and H2.
%-	This explanation implies that the effect of steam would be mainly to prevent the formation of PAHs in the first place, rather than supporting their destruction.

Both effects of steam on tar do require a sufficient temperature, which is the primary factor affecting tar cracking efficiency. Most of the experimental results on tar reforming with steam found in the literature used externally heated reactors, maintained at temperatures between 700 and 900°C to isolate the chemical effect of steam. However, experimental research on autothermal staged downdraft gasifiers with steam have reported a higher tar concentration, due to a negative impact on the reaction temperature \cite{Jarungthammachote2012, DeSales2017}.


\section{Tar sampling and analysis}

Tars are sampled from the hot syngas flow using a tar sampling set-up based on the Tar Protocol guidelines \cite{Neeft2002}, illustrated in Figure \ref{fig:tarprotocol}. Syngas is sampled between the dust cyclone and the condenser, and kept above tar dew point by using trace heating and insulation until a hot filter used to remove particles. Syngas then passes through seven impingers filled with 50ml isopropanol, except for the first and last of them. The first four impingers are located inside a water bath at ambient temperature, while the last three are in a glycolated water bath inside a cooled chamber at -20°C to ensure complete tar condensation. Silica gel is used as a back-up protection of the downstream pump and instruments. The total flow of dry syngas that is sampled is measured by a volumetric meter, and temperature and pressure are continuously recorded for volume correction to Nm$^3$. After tar sampling, \ce{N2} is used to flush the impingers line, then the syngas sampling line. During the experimental campaign, tar sampling durations were between 38 and 84 minutes, for total volumes of dry syngas between 0.11 and 0.29\,Nm$^3$.

After sampling, impingers are emptied and rinsed with isopropanol into a 500ml bulk solution flask. The flask is weighted before and after filling so that the total amount of tar solution is known. The tar solutions are analysed by GCMS with 84 species quantified: the list of tar species is available in Table \ref{tab:tarspecies} in the Appendix. The concentrations are measured in mg/l of isopropanol solution, and converted to mg/Nm$^3$ of dry syngas thanks to the total amount of solution and the total volume of sampled syngas.

\begin{figure}[!ht]
	\centering
	\includegraphics[width=\textwidth]{TarProtocolPID.pdf}
	\caption{Tar sampling line, based on the Tar Protocol Guidelines \cite{Neeft2002}.}
	\label{fig:tarprotocol}
\end{figure}


\section{Experimental results}

Among the 84 species looked for in the GCMS analysis of the tar samples (Table \ref{tab:tarspecies}), only 15 were found in a representative quantity, higher than their Limit Of Quantification (LOQ). The LOQ of most tar species is 1\,mg/l, in the isopropanol solution. The equivalent value, in mg/Nm$^3$ of dry syngas, depends on the conversion factor specific to each sampling conditions. For most runs, the equivalent LOQ of most species is between 1.8 and 2.5\,mg/Nm$^3$ of dry syngas (Table \ref{tab:tarsampling}).

%This accuracy is insufficient to calculate the tar dew point as 1\,mg/Nm$^3$ of typical class V tar species would start to condense at around 150°C \cite{Kiel2004}, thus have a very high impact on the tar dew point even at very low concentrations.

In the reference air conditions, the total concentration of tar species heavier than benzene lies around 30\,mg/Nm$^3$ (Figure \ref{fig:tar_tot_mgNm3}). Adding steam to air does not have any distinguishable effect on the total tar concentration. Replacing air by \ce{O2} increases tar concentration to 45-60\,mg/Nm$^3$. Adding steam to oxygen results in a strong increase of the tar concentration: doubling it to 120\,mg/Nm$^3$ at SOR=2.

Benzene concentration is much higher than all other tar species, which is why it is treated separately. In the reference air conditions, its concentration is in the range 100--200\,mg/Nm$^3$. It increases to 300--400\,mg/Nm$^3$ under oxygen gasification, and up to 450--650\,mg/Nm$^3$ with oxygen and $\mathrm{SOR}=2$. No clear effect of steam injection with air is identified.

\begin{figure}[!ht]
	\centering
	\includegraphics[width=0.7\textwidth]{tar_tot_mgNm3.pdf}
	\caption{Concentration of tars heavier than benzene increases from 30\,mg/Nm$^3$ with air up to 120\,mg/Nm$^3$ in oxy-steam conditions.\newline\textit{Lines: regression model $f(\mathrm{SOR},x_{\ce{O2}},x_{\ce{O2}}\times\mathrm{SOR})$. Markers: experimental points.}\newline\textit{Semi-transparent points are outliers.}}
	\label{fig:tar_tot_mgNm3}
\end{figure}

\begin{figure}[!ht]
	\centering
	\includegraphics[width=0.7\textwidth]{benzene_mgNm3.pdf}
	\caption{Concentration of benzene is 3--5 times higher than other tar species. It is highest under oxy-steam conditions.\newline\textit{Lines: regression model $f(\mathrm{SOR},x_{\ce{O2}},x_{\ce{O2}}\times\mathrm{SOR})$. Markers: experimental points.}\newline\textit{Semi-transparent points are outliers.}}
	\label{fig:benzene_mgNm3}
\end{figure}

For a more detailed understanding of how tar production is affected by varying conditions, tar classes II, III and IV are separately analyzed in Figures \ref{fig:tar_classIV} to \ref{fig:tar_classII}. Observing only the tar concentration in mg/Nm$^3$ can hide effects that are solely related to the variable syngas dilution, with more or less \ce{N2}. To avoid this effect and get more insight into the production of tar, the amount is normalized by the syngas LHV, expressed in mg/MJ. The concentration of tar in the dry syngas is given for each class in Table \ref{tab:tar_mgNm3_full} in the Appendix.

Class IV tars found in the samples are mainly composed of naphthalene (more than $90\%$), with some acenaphthylene, acenaphthene and phenanthrene. They are tertiary species produced by the thermal cracking of secondary tars, and are problematic due to their relatively high dew point, as introduced in Section \ref{sec:theo_tar}.

When added to air, steam decreases the production of tertiary tar: from 2--4\,mg/MJ down to 1\,mg/MJ (Figure \ref{fig:tar_classIV}). This effect is probably due to the presence of more \ce{H^.} and \ce{OH^.} radicals, which prevent the formation of PAHs by saturating aromatic rings and reforming their phenolic precursors. On the other hand, the replacement of air by oxygen slightly increases class IV tar production (5--6\,mg/MJ), despite an increase of more than 200°C of the average temperature in the oxidation zone. When adding steam to oxygen, the class IV tar production further increases to 8--12\,mg/MJ at SOR=1 and up to 15\,mg/MJ at SOR=2. That effect is correlated with a lower $T_\text{oxi,avg}$ resulting from the injection of steam. 

These last trends are unexpected, as steam and oxygen are in theory beneficial to an efficient tar cracking. The main explanation is that the use of oxygen can imply a less efficient mixing between \ce{O2} and the pyrolysis volatiles in the oxidation zone, as the volume flow from the secondary nozzles is 5 times lower than with air. This can leave more areas in very rich conditions at a very high temperature, favorable to the recombination of hydrocarbons in polyaromatic rings. With oxygen only, this effect appears partly offset by the higher average temperature favorable to tar cracking. The addition of steam reduces the average temperature in the oxidation zone (as discussed in Section \ref{sec:res_Toxi}), affecting the efficiency of thermal tar cracking.

%The main theoretical effect of steam is to enhance the conversion of phenolic compounds (class II) into single-ring aromatics and light PAHs. As the amount of phenol is already very low without steam, its presence does not seem to bring much benefit, while still boosting class III and IV tar. Since class V tar (heavy PAHs) could not be quantified due to their too low concentration, we do not know whether steam simultaneously decreased them or not. However, even if they are reduced, the negative effect on class IV tar is probably not compensated.

Class III tar species are essentially benzene (94--100\%), with some toluene (up to 6\%) and small amounts of styrene (max 0.5\%). These species have a very low dew point and are relatively easy to burn, and are thus not harmful for most applications. Yet, they are relevant to monitor as their presence can affect the syngas combustion properties, such as ignition timing and flame radiation.

The production of Class III species is 7--40 times higher than Class IV species, with at least 25\,mg/MJ under air gasification conditions (Figure \ref{fig:tar_classIII}). They are affected in a quite similar way as what is observed for Class IV species, as they are also tertiary products: their production increases slightly when switching from air to oxygen (from 25--45 to 45--60\,mg/MJ), and adding steam worsens that effect by lowering the oxidation zone temperature (up to 100\,mg/MJ). Regarding air-steam conditions, the effect of steam is not clear but it tends to slightly increase the production of benzene; this is an expected effect, as steam favors single-ring aromatics.


\begin{figure}[!ht]
	\centering
	\includegraphics[width=0.7\textwidth]{tar_classIV_mgMJ.pdf}
	\caption{Up to 5 times more class IV tars are produced in oxy-steam conditions. Steam addition to air does reduce class IV tar production.\newline\textit{Semi-transparent markers are outliers.}}
	\label{fig:tar_classIV}
\end{figure}

\begin{figure}[!ht]
	\includegraphics[width=0.7\textwidth]{tar_classIII_mgMJ.pdf}
	\caption{Class III tars are affected by oxygen in a similar way as Class IV species.\newline\textit{Semi-transparent markers are outliers.}}
	\label{fig:tar_classIII}
\end{figure}


Class II tar species are mainly phenol (42\% or more), cresol (up to 32\%) and xylenol (max 15\%) in the present results. These species are problematic on two levels: their high dew point, and their polarity which makes the waste-water treatment very complex, as explained in Section \ref{sec:theo_tar}. They are produced by the conversion of primary pyrolysis products, and most of them should normally be destroyed or converted to tertiary tars in the oxidation zone. Their presence can be due to a too low pyrolysis front, resulting in incompletely pyrolysed biomass falling to the reduction zone, some secondary species thereby avoiding the oxidation zone.

In all experiments with \ce{O2} or air-\ce{O2}, the amount of class II tars is below 1\,mg/MJ. They are only found in substantial quantity in some results of the first experimental campaign, studying the impact of steam addition with air. The highest amount (14\,mg/MJ) is obtained without steam injection, down to 6\,mg/MJ at SOR=2. This could indicate a positive effect of steam, favoring their reforming by providing many \ce{H^.} and \ce{OH^.} radicals. However, the lowest amounts are also associated with a slightly higher and warmer pyrolysis front, which could be another explanation to the lower production of Class II tars. Overall, their amount is very low in almost all results, and no clear impact of steam and oxygen can be concluded.%On the other hand, all results of the second experimental campaign exhibit very low class II tar production, including points for which the average pyrolysis zone temperature is even lower.

\begin{figure}[!ht]
	\centering
	\includegraphics[width=0.7\textwidth]{tar_classII_mgMJ.pdf}
	\caption{Class II tar species were only found in substantial amount during the "air-steam" campaign. Their amount is lower with steam injection, which also corresponds to a slightly higher and warmer pyrolysis front (+30°C average temperature).}
	\label{fig:tar_classII}
\end{figure}

%\begin{figure}[!ht]
%	\centering
%	\includegraphics[width=0.7\textwidth]{Tpyro_example.pdf}
%	\caption{Temperature distribution within the pyrolysis zone, for four experimental conditions. Left-to-right values correspond to equidistant thermocouples (11\,cm separation) from the top to the bottom of the pyrolysis zone. Boxes indicate the distribution during each run. The average temperature is shown as grey dots.}
%	\label{fig:Tpyro_example}
%\end{figure}
